%Related work section

\subsection{Social network effects upon financial markets}
%With social networks playing an ever increasing role in peoples life, all kinds of discussions take place on forums such as Reddit or Twitter. 
The impact of online social networks on society and financial markets was intensively studied from various angles. Especially Twitter received notable attention \citep{lehrer2021does, Gu2020,Behrendt2018, Yang2015}, but also StockTwits \citep{Audrino2020, chen2022sonic} and sentiment from across sources \parencite{audrino2024does, Gan2020, Jiao2020}. Often network activity such as aggregated number of posts was found to be an important factor in modelling network influence on markets. The actions and discussions on the subreddit WSB on Reddit, which led to the GME market frenzy, resulted in yet another social network impacting what happens in the financial markets \citep{Pedersen2022,Klein2022,Umar2021}. 
%The effect of the market frenzy was so detrimental, that even a US senate hearing took place, during which they questioned the main actors involved in the market frenzy, including one of the Reddit users hyping GME. 
In the aftermath of the GME frenzy, it became obvious that the `internet crowd' represents a new force driving financial markets, see e.g. the study by \textcite{Pedersen2022} or \textcite{Yang2015}. Indeed, attention within a community is well known to impact stock prices \citep{Andrei2014}. The `meme stock phenomenon' \citep{Securities2021} leads regularly to herding behaviour and market excesses, impacting asset prices \citep{Aloosh2021, Costola2021}. Our study adds to this literature by using a large Reddit data set and modelling the Reddit network, determining when and how WSB users post and how that relates to the GME stock price. The model we develop is illustrated on the GME network but has the potential to be applied to any other meme stock network. 


\subsection{Changes in social networks and their link to financial markets}
Network theory states that participants in a social network update their beliefs by listening to others in a social network \citep{DeGroot1974}. Whereas the theory is general for any kind of social network, the emergence of online social networks expanded its applications. If the social network is about stocks such as the subreddit Wallstreetbets, then the users update their beliefs about the value of a stock via the social network. In the model of \citet{Pedersen2022}, the changes in social network users' beliefs about an asset determine its price on the markets. Hence, the changes in users beliefs about an asset would lead to buying or selling the asset until the price matches users beliefs again. 
%Hence, if the users change their beliefs about the price of an asset and as by their updated belief the asset is trading for a different price, then this would lead to buying or selling the asset until the price matches users beliefs again. 
%The DeGroot model states that users share their opinion with a certain probability with others, resulting in the updating of their opinions, which motivates the modeling of interactions on social networks with exponential random graph models \citep{Deng2024, Bhattacharya2019}. %The co-occurrence of changes in beliefs and behaviours in the social network has been observed by \textcite{Bailey2024}. 
Our study contributes to this literature by including a change point mechanism to model belief updates on the individual node behaviour. Our developed model then links these changes in behaviour to the stock market which enables us to empirically quantify the effect of meme stock social networks on the market.


\subsection{Link formation models and the role of users}
Link formation in social networks is a widely studied topic. Random network models such as Exponential Random Graph Models and Latent Variable Models have been widely studied for such tasks \citep{kim2018review,Deng2024}. Not all of these models are applicable for our purpose, as, for example, many become intractable, which makes them (only) suitable for small networks \citep{Yon2021}. Recently, there have been studies using deep learning models for link prediction \citep{barros2021survey,yu2023towards,beladev2023graphert,hui2024domain}. However, these models are not easily interpretable, and therefore provide little insight into the drivers of individuals posting behaviour. For large networks another option is to use logistic regression-based dynamic network models \citep{agaskar2013alarm,almquist2013dynamic}, because these provide a scalable framework for inference on large-scale temporal networks \citep{mallik2019stable}. In addition, these models naturally allow for incorporation of network dynamics through the autoregressive modelling structures. Furthermore, \textcite{bhattacharya2019coevolution} found that variables such as the past posting frequency and a nodes' out-degree matter for their future communication behaviour, opening the pathway to predict the probability of communication between two users based on their past behaviour and relevance to the network. This motivates our adoption of an individual-level logistic regression type model which is scalable as well as interpretable and allows naturally for incorporation of network dynamics. %specific model , containing variables representing the involved users own past behaviour and the ones of its peers, allowing for an explainable framework for communication behaviour. 

Link formation in social networks is strongly interlinked with the role users have within networks. Naturally users `trust' certain other users more, which can lead to the emergence of opinion leaders \citep{Lu2013}, also termed influencers \citep{han2020importance, chen2022sonic}. Influencers are typically characterized by their topological position in the network, for example as measured by various centrality measures \citep{gelper2021competition}. The difference in user relevance on the social network is known to explain user behaviour on a social network \citep{Lu2013,zhang2018direct} as well as its own behaviour and the one of its peers. It has also been shown that influencers do influence one another \citep{Ji2025}. Our study builds on these findings by modelling separate models for vital and non-vital users in the network (one may think of this terminology as influencers and non-influencers), enabling us to study the effect of these different roles in the network on the markets.

%Taking into account that nodes within a network are heterogeneous, implies that a different set of determinants impacts users' communication behaviour with another node. Communication between users on online social networks is also known for being irregular and infrequent, causing the observed network to be sparse. Consequently, a model is needed that contains a large set of candidate variables to explain nodes' communication behaviour while considering that not all variables matter for each link, motivating the introduction of sparse link prediction. 

\subsection{Change point detection}
The detection of changes in a system has been a long-standing research topic. Changes are typically identified by comparing models before and after a possible change point. The overarching aim of change point detection is to identify a point in time where the network or a node in the network exhibits a change in behaviour \citet{ghezzi2024fast}. Within the field of network analysis, there is an extensive literature \citep{kendrick2018change, woodall2017overview}.  
Most of this literature considers networks with a small number of nodes, typically no more than a few thousand \citep{chen2024monitoring,Yu2022monitoring}. In our experience, the huge size of social networks in combination with their typical sparsity often challenges the existing algorithms for change point detection due to their computational complexity. In addition, most literature focuses on so-called global change point detection, identifying time points where the whole network, or a substantial part of the network, changes. An example is \citet{zwetsloot2025online} who identify change points in the same GME data, however, these are at a network community level, providing no-change identification at the nodal level. In contrast, the target of this study requires change point detection on the nodal level. Because, it is known that nodes have heterogeneous effects on the network \citet{chen2023community}, especially highly influential nodes. Also, as networks are highly dynamic, users may change their behaviour multiple times within our observation period, hence we need algorithms for multiple change point detection \citep{niu2016multiple,bai2023multiple}. This motivates our adoption of multiple change points detection at the node level.

%MOVED TO SEC 4 For our proposed method, we tackle the network size challenge by estimating a dedicated model for each pair of users. This allows for the analysis of a network of huge nodal size (500000+ in our case) while relying upon the vast literature for change point detection in time series. It also enables us to look for changes at the nodal level. To keep the identification of change points flexible, greedy algorithms such as in \textcite{bai2023multiple} allow data-driven detection of the change points. 
%discuss De Groot and other change detection

\subsection{High-dimensional networks and sparsity}
Focusing upon estimating the communication probability between any two users is a popular approach allowing to analyse the microstructure of a network (for example, \citet{agaskar2013alarm}). This implies that a tailored model per user is needed to understand the drivers behind their behaviour. Not all parameters included in the model may be relevant for every users' behaviour, hence the system may experience sparsity. In fact, many networks are characterised by huge dimensionality, but have sparse links \citep{paula_econometrics_2017}. The high-dimensional nature of networks challenges the joint analysis of nodes and identification of important links. Methods from sparse modelling, originally developed for linear models \citep{tibshirani_regression_1996, fan_variable_2001, zou_adaptive_2006}, were increasingly used for network analysis to handle the sparse nature of the underlying models \citep{Barigozzi2019, Trimborn2022a, Trimborn2022, Zhang2023}. The issue of sparsity within social networks concerns various modelling challenges, ranging from link prediction \citep{Trimborn2022} to change point detection \citep{chen2024monitoring}. Likewise, the network model developed in this study incorporates a LASSO type estimator \citep{tibshirani_regression_1996} to account for the observed sparse structure. 
